\documentclass[11pt]{article}
\usepackage[T1]{fontenc}
\usepackage{lmodern}
\usepackage{amsmath,amssymb}
\usepackage[a4paper,margin=28mm]{geometry}
\usepackage{microtype}
\usepackage{xcolor}
\usepackage{enumitem}
\usepackage{listings}
\usepackage[hidelinks]{hyperref}
\usepackage[nameinlink,noabbrev]{cleveref}
\setlength{\emergencystretch}{3em}
\definecolor{codegray}{RGB}{245,245,245}
\newcommand{\workflowstep}[1]{\subsection{Workflow Step #1}}
\newcommand{\taupar}{\tau}
\newcommand{\taumod}{\tau_{\mathrm{mod}}}
\newcommand{\EJ}{E_J}
\newcommand{\RamF}{f_{\mathrm R}}
\newcommand{\TODO}[1]{\par\noindent\textcolor{red!70!black}{\textbf{TODO:} #1}\par}
\lstdefinestyle{parigp}{language={},basicstyle=\ttfamily\small,
  backgroundcolor=\color{codegray},frame=single,breaklines=true,
  columns=fullflexible,keepspaces=true,showstringspaces=false}

\title{From Six-Parameter Somos--4 Data to Weierstrass, Jacobi,\\
Ramanujan Theta Representations\\[0.4em]
\large A Constructive Workflow}
\author{Thomas Scheuerle}
\date{}

\begin{document}
\maketitle

\begin{abstract}
A generic Somos--4 orbit with four freely prescribed initial values satisfies
\[
 H_{n+4}H_n=\alpha H_{n+3}H_{n+1}+\taupar H_{n+2}^2.
\]
An earlier workflow treated the normalized case $\alpha=1$ and organized the
passage from a direct algebraic generating function to a genus-one curve, a
Weierstrass sigma representation, and a Jacobi $\theta_1$ representation.
The present paper develops the corresponding six-parameter workflow and adds
a further endpoint suggested by Michael Somos: Ramanujan's general theta
function
\[
 \RamF(a,b)=\sum_{k\in\mathbb Z}
 a^{k(k+1)/2}b^{k(k-1)/2},\qquad |ab|<1.
\]
Starting from $(\lambda,m,r,\eta,\alpha,\taupar)$, we separate three layers that
must not be conflated: the Somos orbit, its elliptic curve, and a chosen
power-series lift.  A weighted transport from the five-parameter algebraic
generating function introduces a sixth root of $\alpha$ in that particular
lift, whereas the recurrence and a birational quartic model remain rational in
the six parameters.  The elliptic part of the workflow then proceeds through
periods, a translation point, Abel--Jacobi coordinates, and the classical
sigma formula.  The Jacobi endpoint is converted exactly into Ramanujan form,
with moving arguments $a_n,b_n$ satisfying the fixed-product relation
$a_nb_n=q^2$, where $q=e^{\pi i\taumod}$.  This exposes both a bilateral
$q$-series and an infinite-product representation of the same Somos--4 orbit.
The emphasis is on normalization contracts, branch and indexing checks, and a
reproducible computational plan.  PARI/GP routines and numerical tables are
left as explicit placeholders for a companion implementation.
\end{abstract}

\section{Introduction}
The sigma-function description of generic Somos--4 sequences is classical.
The practical task, however, is not merely to quote the existence of such a
description.  Starting from concrete recurrence data, one must identify the
correct invariant curve, choose compatible coordinates and periods, locate the
translation point, fix Abel--Jacobi lifts, and recover the elementary gauge
factors.  Each step carries a convention whose unnoticed change may preserve
some low-order checks while invalidating the final formula.

The preceding workflow paper carried out this program for the normalized
recurrence
\[
 H_{n+4}H_n=H_{n+3}H_{n+1}+\taupar H_{n+2}^2,
\]
starting from the direct algebraic generating function with prescribed
$(H_1,H_2,H_3,H_4)=(\lambda,m,r,\eta)$.  Two developments motivate the present
follow-up.  First, the generating-function construction has been extended to
the six-parameter recurrence with a free coefficient $\alpha$.  Second,
Michael Somos pointed out that Ramanujan theta functions provide a natural and
less frequently documented endpoint of the workflow.  The latter is not only
a change of notation: Ramanujan's $\RamF(a,b)$ simultaneously displays the
bilateral series and Jacobi triple product, while the sequence index enters
explicitly through its two arguments.

The resulting workflow is
\[
(\lambda,m,r,\eta,\alpha,\taupar)
\longmapsto (J,\mathcal C)
\longmapsto E
\longmapsto P
\longmapsto(\omega_1,\omega_2;z_0,\delta)
\longmapsto
\begin{cases}
 \sigma,\\
 \theta_1,\\
 \RamF(a,b).
\end{cases}
\]
The six-parameter entry requires care.  The direct transported generating
series lives naturally over a radical extension $q_6^6=\alpha$, but the
ordinary Hankel orbit and a rational quartic model live over
\[
 K=\mathbb Q(\lambda,m,r,\eta,\alpha,\taupar).
\]
Accordingly, this paper distinguishes statements about a generating series
from statements about the curve or recurrence.

\subsection{Scope and contribution}
This paper provides:
\begin{itemize}[leftmargin=2em]
 \item a workflow-level synthesis of the six-parameter generating-function
       update with the algebraic-to-sigma construction;
 \item a rational quartic entry point for the six-parameter orbit, together
       with explicit checkpoints separating it from a radical power-series
       lift;
 \item the exact conversion from the Jacobi $\theta_1$ endpoint to
       Ramanujan's general theta function $\RamF(a,b)$;
 \item the fixed-product law $a_nb_n=q^2$ and the associated bilateral-series
       and infinite-product forms;
 \item a verification architecture for recurrence, curve, sigma, Jacobi, and
       Ramanujan evaluations; and
 \item marked placeholders for the later PARI/GP companion and numerical
       examples.
\end{itemize}
It does not claim a new proof of the classical sigma representation, nor does
it assert that every six-parameter algebraic lift requires radicals.  The
non-descent statement concerns the particular weighted transport used below.

\section{Six-parameter input and the three mathematical layers}
Let
\begin{equation}
 H_{n+4}H_n=\alpha H_{n+3}H_{n+1}+\taupar H_{n+2}^2,
 \qquad
 (H_1,H_2,H_3,H_4)=(\lambda,m,r,\eta).
 \label{eq:somos-six}
\end{equation}
We assume throughout the main workflow that all displayed denominators are
nonzero and that the resulting genus-one curve is nonsingular.  Degenerate
cases require separate projective or limiting treatments.

Three related objects must be kept distinct:
\[
 \text{ordinary Hankel orbit},\qquad
 \text{genus-one curve},\qquad
 \text{chosen generating series}.
\]
The recurrence generates $H_n\in K$.  The six-parameter follow-up also gives a
quartic model over $K$.  By contrast, the weighted algebraic-series transport
uses $q_6$ with $q_6^6=\alpha$ and generically has coefficients outside $K$.
This distinction is essential for a correct workflow.

\subsection{Weighted generating-series transport}
Let $G_{1;\lambda,m,r,\eta,\taupar}(x)$ denote the five-parameter series.  Put
\begin{equation}
 \lambda'=\lambda,\quad m'=q_6^{-2}m,\quad r'=q_6^{-6}r,
 \quad \eta'=q_6^{-12}\eta,\quad \taupar'=q_6^{-8}\taupar,
 \qquad q_6^6=\alpha,
 \label{eq:weighted-parameters}
\end{equation}
and define
\begin{equation}
 G_{\alpha;\lambda,m,r,\eta,\taupar}(x)
 =G_{1;\lambda',m',r',\eta',\taupar'}(q_6x).
 \label{eq:weighted-series}
\end{equation}
If $b_k=q_6^ka_k$, row and column factorization gives
\begin{equation}
 H_N(b)=q_6^{N(N-1)}H_N(a),\qquad
 D_N(b)=q_6^{N^2}D_N(a).
 \label{eq:hankel-weights}
\end{equation}
Consequently, the first four ordinary determinants remain
$(\lambda,m,r,\eta)$ and satisfy \cref{eq:somos-six}.  The inherited companion
law is graded:
\begin{equation}
 D_N(b)=q_6^{-3N-2}H_{N+2}(b).
 \label{eq:graded-companion}
\end{equation}
Already $b_1=q_6^{-5}r$, so this transported branch does not generically
descend to $K$.  Nothing in this observation prevents the orbit and curve from
descending.

\section{Reduced dynamics and rational curve model}
Define the gauge-invariant variables
\begin{equation}
 x_n=\frac{H_{n-1}H_{n+1}}{H_n^2}.
 \label{eq:xn}
\end{equation}
Equation \cref{eq:somos-six} implies
\begin{equation}
 x_{n+2}x_nx_{n+1}^2=\alpha x_{n+1}+\taupar,
 \qquad
 x_{n+2}=\frac{\alpha x_{n+1}+\taupar}{x_nx_{n+1}^2}.
 \label{eq:qrt-six}
\end{equation}
The natural first integral is
\begin{equation}
 J=x_nx_{n+1}+\alpha\left(\frac1{x_n}+\frac1{x_{n+1}}\right)
   +\frac{\taupar}{x_nx_{n+1}}.
 \label{eq:J-six}
\end{equation}
For the initial determinants,
\begin{equation}
 x_2=\frac{\lambda r}{m^2},\qquad
 x_3=\frac{m\eta}{r^2},
 \label{eq:x2x3}
\end{equation}
so that
\begin{equation}
 J=\frac{\lambda\eta}{mr}
 +\alpha\frac{m^2}{\lambda r}
 +\alpha\frac{r^2}{m\eta}
 +\frac{\taupar mr}{\lambda\eta}.
 \label{eq:J-seeds}
\end{equation}
A rational quartic model compatible with the six-parameter follow-up is
\begin{equation}
 \mathcal C_{J,\alpha,\taupar}:\qquad
 y^2=(X^2-JX+\taupar)^2-4\alpha^2X.
 \label{eq:quartic-six}
\end{equation}
For $\alpha=1$, this reduces to the quartic used in the preceding workflow.

\TODO{In the code session, verify \cref{eq:J-six,eq:quartic-six} symbolically
from consecutive QRT states and record the exact state-to-quartic map.  The
normalization must be checked rather than inferred solely by specialization.}

\section{The constructive workflow}

\workflowstep{1: Validate the six-parameter input}
Read $(\lambda,m,r,\eta,\alpha,\taupar)$, generate several exact recurrence
values, and test every denominator required by the selected affine charts.
Compute $x_2,x_3$ and $J$ from \cref{eq:x2x3,eq:J-seeds}.

\paragraph{Checkpoint.}
Verify \cref{eq:J-six} for several consecutive pairs generated by
\cref{eq:qrt-six}.  Verify independently that the $H_n$ and $x_n$ routes agree.

\workflowstep{2: Select the orbit, curve, or series route}
If only the elliptic or transcendental representation is required, work over
$K$ with \cref{eq:quartic-six}.  If the direct algebraic generating function is
also required, construct \cref{eq:weighted-series} over $K(q_6)$ and retain the
graded companion law \cref{eq:graded-companion}.  Do not use the literal
five-parameter companion identity after transport.

\paragraph{Checkpoint.}
Confirm $(H_1,H_2,H_3,H_4)=(\lambda,m,r,\eta)$ and
\[
 H_5=\frac{\alpha\eta m+\taupar r^2}{\lambda},\qquad
 H_6=\frac{\alpha H_5r+\taupar\eta^2}{m}.
\]
For the transported series, confirm $D_N=q_6^{-3N-2}H_{N+2}$ at the first
orders.

\workflowstep{3: Convert the rational quartic to Weierstrass form}
Choose one point at infinity as origin and compute a birational map
\begin{equation}
 (X,y)\longleftrightarrow(u,v),\qquad
 E:\ v^2=4u^3-g_2u-g_3.
 \label{eq:weierstrass-general}
\end{equation}
The formulas must specialize to those of the normalized workflow at
$\alpha=1$.

\TODO{Insert the exact six-parameter birational maps, $g_2(J,\alpha,\taupar)$,
$g_3(J,\alpha,\taupar)$, the elliptic discriminant, and the affine exceptional
set after symbolic derivation and PARI/GP verification.}

\paragraph{Checkpoint.}
Verify both map compositions modulo the defining curve equations.  Check that
the quartic discriminant and $g_2^3-27g_3^2$ agree up to the stated scaling,
and verify the $j$-invariant independently.

\workflowstep{4: Identify the fixed elliptic translation}
Map the QRT state $(x_n,x_{n+1})$ to a point $Q_n\in E$.  Determine the fixed
point $P\in E$ satisfying
\begin{equation}
 Q_{n+1}=Q_n+P.
 \label{eq:translation}
\end{equation}
The sign convention for the second quartic coordinate fixes whether a step is
represented by $P$ or $-P$.

\TODO{Derive and insert the explicit six-parameter state map and translation
point.  Verify that the formulas reduce to
$P=((J^2-4\taupar)/12,1)$ when $\alpha=1$ in the text convention.}

\paragraph{Checkpoint.}
For several exact QRT iterates, verify that $Q_n$ lies on $E$ and that
$Q_{n+1}-Q_n=P$ is independent of $n$.

\workflowstep{5: Compute and normalize the period lattice}
Compute an oriented basis of full periods $(2\omega_1,2\omega_2)$ for
\cref{eq:weierstrass-general} and set
\begin{equation}
 \taumod=\frac{\omega_2}{\omega_1},\qquad
 \operatorname{Im}(\taumod)>0.
 \label{eq:tau-mod}
\end{equation}
Only the lattice is intrinsic; the ordered basis is a retained workflow
choice.  Every later sigma, theta, and Ramanujan evaluation must use the same
basis.

\paragraph{Checkpoint.}
Reconstruct $g_2,g_3$ from the lattice and compare with the algebraic values.
State explicitly whether software returns full or half periods.

\workflowstep{6: Compute Abel--Jacobi coordinates}
Choose lifts $z_2$ and $\delta$ such that
\[
 (\wp(z_2),\wp'(z_2))=Q_2,
 \qquad
 (\wp(\delta),\wp'(\delta))=P,
\]
and define, for the present $H_1$-based indexing,
\begin{equation}
 z_0=z_2-3\delta\pmod\Lambda.
 \label{eq:z0}
\end{equation}
Then $Q_n$ is represented by $z_0+(n+1)\delta$.

\paragraph{Checkpoint.}
Compare both Weierstrass coordinates, not only the first coordinate, for
several $n$.  This fixes the sign of $\delta$.

\workflowstep{7: Fix the sigma normalization}
Use the classical representation
\begin{equation}
 H_n=A B^n
 \frac{\sigma(z_0+n\delta)}{\sigma(\delta)^{n^2}}.
 \label{eq:sigma}
\end{equation}
Put
\[
 S_n=\frac{\sigma(z_0+n\delta)}{\sigma(\delta)^{n^2}}.
\]
When the displayed factors are nonzero, the first two initial values give
\begin{equation}
 B=\frac{mS_1}{\lambda S_2},\qquad
 A=\frac{\lambda^2S_2}{mS_1^2}.
 \label{eq:AB}
\end{equation}
The values $H_3=r$ and $H_4=\eta$ are then independent consistency checks.
The coefficient $\alpha$ is encoded in the elliptic data and the translation
step, not in an extra arbitrary sigma gauge.

\paragraph{Checkpoint.}
Compare \cref{eq:sigma} with exact recurrence values and verify
\begin{equation}
 x_n=\wp(\delta)-\wp(z_0+n\delta)
 \label{eq:x-sigma}
\end{equation}
under the selected normalization.  If the six-parameter scaling changes this
identity, record the corrected factor explicitly rather than absorbing it
silently.

\workflowstep{8: Convert sigma to Jacobi $\theta_1$}
For the chosen ordered half periods, let
\[
 \eta_1=\zeta(\omega_1),\qquad
 c=\frac{2\omega_1}{\pi\theta_1'(0\mid\taumod)}.
\]
The normalization contract is
\begin{equation}
 \sigma(z)=
 \exp\!\left(\frac{\eta_1z^2}{2\omega_1}\right)c\,
 \theta_1\!\left(\frac{\pi z}{2\omega_1}\,\middle|\,\taumod\right).
 \label{eq:sigma-theta}
\end{equation}
Define
\begin{equation}
 A_0=\frac{\pi\delta}{2\omega_1},\qquad
 B_0=\frac{\pi z_0}{2\omega_1}.
 \label{eq:A0B0}
\end{equation}
Then
\begin{equation}
 H_n=K_0q_0^{n^2}r_0^n
 \theta_1(A_0n+B_0\mid\taumod),
 \label{eq:theta}
\end{equation}
where
\begin{align}
 K_0&=Ac\exp\!\left(\frac{\eta_1z_0^2}{2\omega_1}\right),\\
 q_0&=\exp\!\left(\frac{\eta_1\delta^2}{2\omega_1}\right)
       \sigma(\delta)^{-1}
     =\bigl(c\theta_1(A_0\mid\taumod)\bigr)^{-1},\\
 r_0&=B\exp\!\left(\frac{\eta_1z_0\delta}{\omega_1}\right).
 \label{eq:theta-constants}
\end{align}

\paragraph{Checkpoint.}
Evaluate the sigma and Jacobi forms through independent code paths.  Libraries
that use the nome rather than $\taumod$, or a rescaled first argument, require
an explicit adapter.

\workflowstep{9: Convert Jacobi $\theta_1$ to Ramanujan $\RamF(a,b)$}
Set
\begin{equation}
 q=e^{\pi i\taumod},\qquad |q|<1,
 \label{eq:nome}
\end{equation}
and define Ramanujan's general theta function by
\begin{equation}
 \RamF(a,b)=\sum_{k=-\infty}^{\infty}
 a^{k(k+1)/2}b^{k(k-1)/2},\qquad |ab|<1.
 \label{eq:ramanujan-f}
\end{equation}
Its Jacobi triple product is
\begin{equation}
 \RamF(a,b)=(-a;ab)_\infty(-b;ab)_\infty(ab;ab)_\infty.
 \label{eq:triple-product}
\end{equation}
For the convention in \cref{eq:sigma-theta}, the exact bridge is
\begin{equation}
 \theta_1(z\mid\taumod)
 =i q^{1/4}e^{-iz}
 \RamF\!\left(-e^{2iz},-q^2e^{-2iz}\right).
 \label{eq:theta-ramanujan}
\end{equation}
Put $z_n=A_0n+B_0$ and define
\begin{align}
 a_n&=-e^{2iz_n},&
 b_n&=-q^2e^{-2iz_n},\\
 K_R&=iK_0q^{1/4}e^{-iB_0},&
 r_R&=r_0e^{-iA_0}.
 \label{eq:ramanujan-data}
\end{align}
Substitution into \cref{eq:theta} gives the Ramanujan endpoint
\begin{equation}
 \boxed{
 H_n=K_Rq_0^{n^2}r_R^n\RamF(a_n,b_n).}
 \label{eq:ramanujan-somos}
\end{equation}
The moving arguments satisfy the fixed-product law
\begin{equation}
 \boxed{a_nb_n=q^2,}
 \label{eq:fixed-product}
\end{equation}
which is independent of the Somos index $n$.

\paragraph{Checkpoint.}
Verify \cref{eq:theta-ramanujan} first for generic complex test arguments, then
compare \cref{eq:ramanujan-somos} with the Jacobi, sigma, and recurrence values.
Test convergence as the bilateral truncation grows.

\section{What the Ramanujan endpoint adds}
The conversion is exact, but it reveals structure that is less visible in the
sigma notation.  From \cref{eq:ramanujan-f,eq:ramanujan-data},
\begin{equation}
 \RamF(a_n,b_n)=
 \sum_{k\in\mathbb Z}(-1)^kq^{k(k-1)}e^{2ik(A_0n+B_0)}.
 \label{eq:bilateral}
\end{equation}
Thus
\begin{equation}
 H_n=K_Rq_0^{n^2}r_R^n
 \sum_{k\in\mathbb Z}(-1)^kq^{k(k-1)}
 e^{2ikB_0}e^{2ikA_0n}.
 \label{eq:fourier-index}
\end{equation}
The sequence index occurs explicitly in the exponential factor
$e^{2ikA_0n}$.  It is therefore natural to regard \cref{eq:fourier-index} as a
bilateral theta or Fourier expansion in the discrete index, without assigning
it an interpretation beyond the exact identity.

The product side is equally direct.  Since $a_nb_n=q^2$,
\begin{equation}
 \RamF(a_n,b_n)=
 (e^{2iz_n};q^2)_\infty
 (q^2e^{-2iz_n};q^2)_\infty
 (q^2;q^2)_\infty.
 \label{eq:ramanujan-product}
\end{equation}
Hence the same Somos orbit admits Weierstrass, Jacobi, bilateral $q$-series,
and infinite-product descriptions.

\subsection{Special Ramanujan functions}
Ramanujan's classical specializations include
\[
 \phi(q)=\RamF(q,q),\qquad
 \psi(q)=\RamF(q,q^3),\qquad
 f(-q)=\RamF(-q,-q^2).
\]
The generic arguments in \cref{eq:ramanujan-data} do not usually reduce to
these special forms.  Such reductions require additional arithmetic relations
among $q$, $A_0$, and $B_0$.  They should therefore be treated as special
modular cases rather than as part of the generic workflow.

\TODO{Investigate, after the generic implementation is verified, whether the
classical Somos--4 curve or other arithmetic examples yield useful
$\phi$, $\psi$, eta-product, or modular-equation specializations.  Do not claim
such a reduction without an exact parameter relation.}

\section{Algorithm summary}
\begin{enumerate}[label=\textbf{Step \arabic*.},leftmargin=2.7em]
 \item Read $(\lambda,m,r,\eta,\alpha,\taupar)$ and generate exact recurrence
       data.
 \item Compute $x_2,x_3$ and the invariant $J$.
 \item Decide whether the calculation needs the radical generating-series
       lift or only the rational orbit and curve.
 \item Construct the rational quartic and test nonsingularity.
 \item Compute and verify a birational Weierstrass model.
 \item Identify the fixed translation point from the QRT dynamics.
 \item Compute and orient the period lattice.
 \item Compute $z_0$ and $\delta$, retaining signs and period representatives.
 \item Determine $A,B$ and verify all four initial values.
 \item Convert sigma to Jacobi form using one stated normalization contract.
 \item Convert Jacobi form to Ramanujan $\RamF(a,b)$ and verify
       $a_nb_n=q^2$.
 \item Compare recurrence, sigma, Jacobi, Ramanujan-series, and
       Ramanujan-product evaluations.
\end{enumerate}

\section{Computational architecture}
The companion implementation should keep exact and numerical stages separate.
The proposed components are:
\begin{itemize}[leftmargin=2em]
 \item exact Somos recurrence and QRT iteration for $(\alpha,\taupar)$;
 \item exact invariant and quartic construction;
 \item exact forward and inverse birational maps;
 \item exact curve discriminant and $j$-invariant checks;
 \item numerical periods and elliptic logarithms at controlled precision;
 \item sigma evaluation using the selected lattice;
 \item an independent Jacobi $\theta_1$ implementation or adapter;
 \item direct bilateral summation of $\RamF(a,b)$;
 \item optional evaluation by the triple product; and
 \item a consolidated error report against exact recurrence values.
\end{itemize}

\subsection{PARI/GP placeholders}
\begin{lstlisting}[style=parigp]
\\ TODO: six-parameter recurrence and invariant
somos6_step(v, alpha, tau) = { /* placeholder */ };
qrt6_step(x, y, alpha, tau) = { /* placeholder */ };
qrt6_invariant(x, y, alpha, tau) = { /* placeholder */ };

\\ TODO: rational quartic and verified Weierstrass maps
quartic6_poly(J, alpha, tau, X) =
  (X^2 - J*X + tau)^2 - 4*alpha^2*X;
quartic6_to_weierstrass(X, Y, J, alpha, tau) =
  { /* placeholder */ };
weierstrass_to_quartic6(u, v, J, alpha, tau) =
  { /* placeholder */ };

\\ Ramanujan's general theta function by bilateral truncation
ramanujan_f(a, b, N = 40) =
{
  sum(k = -N, N,
      a^(k*(k+1)/2) * b^(k*(k-1)/2));
};

\\ theta_1 via the exact Ramanujan conversion used in the text
ramanujan_theta1(tau_mod, z, N = 40) =
{
  my(q = exp(Pi*I*tau_mod));
  I*q^(1/4)*exp(-I*z)
    * ramanujan_f(-exp(2*I*z),
                  -q^2*exp(-2*I*z), N);
};

\\ TODO: compare recurrence, sigma, Jacobi, Ramanujan sum and product
workflow6_demo(lambda, m, r, eta, alpha, tau, N = 10, prec = 80) =
  { /* placeholder */ };
\end{lstlisting}

\TODO{In the dedicated code session, replace every placeholder, add adaptive
bilateral truncation and product truncation, and report numerical errors as a
function of precision and truncation order.}

\section{Numerical examples}
\subsection{Normalized classical Somos--4}
For
\[
 (\lambda,m,r,\eta,\alpha,\taupar)=(1,1,1,1,1,1),
\]
the workflow must reproduce
\[
 1,1,1,1,2,3,7,23,59,\ldots.
\]
This example also provides a regression test against the earlier workflow.

\TODO{Insert the recomputed periods, Abel--Jacobi data, sigma constants,
Jacobi constants, nome $q$, several $(a_n,b_n)$ pairs, and a convergence table
for the Ramanujan sum and product.  Recompute rather than copying rounded
values from the earlier manuscript.}

\subsection{A genuinely six-parameter example}
Choose rational nonzero input with $\alpha\ne1$ for which the quartic is
nonsingular and the first several recurrence denominators do not vanish.
The example should show separately:
\begin{enumerate}[leftmargin=2.3em]
 \item exact recurrence and QRT checks over $K$;
 \item the rational quartic and Weierstrass invariants;
 \item the radical coefficient $b_1=q_6^{-5}r$ of the transported series;
 \item numerical sigma and Jacobi evaluations; and
 \item independent Ramanujan-series and product evaluations.
\end{enumerate}

\TODO{Select and verify the example during the PARI/GP session.}

\section{Reliability and reproducibility}
The principal failure modes are now broader than in the normalized workflow:
\begin{itemize}[leftmargin=2em]
 \item confusing the sixth root $q_6$ used in the generating-series transport
       with the elliptic nome $q=e^{\pi i\taumod}$;
 \item applying the ungraded companion identity after weighted transport;
 \item silently changing the quartic or Weierstrass scaling when $\alpha$ is
       introduced;
 \item reversing the QRT translation by changing the sign of a second
       coordinate;
 \item mixing full and half periods;
 \item using inconsistent Abel--Jacobi representatives;
 \item confusing $q$ with $q^2$ in the Ramanujan arguments; and
 \item comparing two formulas that were derived through the same code path
       and therefore do not constitute independent verification.
\end{itemize}
The notation $q_6$ is reserved for $q_6^6=\alpha$, while $q$ denotes the nome.
This typographical separation should be maintained in both manuscript and
code.

\section{Discussion and future directions}
The six-parameter update sharpens the separation between algebraic and
transcendental data.  The ordinary Hankel orbit and its rational quartic may be
handled over $K$, while a distinguished generating-series lift can retain a
radical orientation invisible to ordinary Hankel determinants.  Once the
curve and translation are fixed, the sigma and Jacobi stages proceed as in the
normalized case, but with modified elliptic invariants.

The Ramanujan endpoint adds a complementary view.  It packages the same theta
value as a bilateral series and an infinite product, with a constant nome
product $a_nb_n=q^2$.  This may be particularly useful when the elliptic data
have arithmetic or modular specializations.  The generic identity should be
established and numerically checked first; special identities involving
$\phi$, $\psi$, $f(-q)$, eta products, or modular equations belong to a later
arithmetic analysis.

\section{Conclusion}
The extended workflow begins with six freely varying parameters, preserves the
rational Somos orbit and elliptic curve, and tracks explicitly where a radical
enters the transported algebraic generating series.  It then carries the orbit
through QRT dynamics, a Weierstrass model, periods, Abel--Jacobi coordinates,
and the classical sigma representation.  The final conversion
\[
 \sigma\longleftrightarrow\theta_1\longleftrightarrow\RamF(a,b)
\]
is exact under one stated normalization contract.  In Ramanujan form, the
sequence is
\[
 H_n=K_Rq_0^{n^2}r_R^n\RamF(a_n,b_n),
 \qquad a_nb_n=q^2,
\]
which exposes both the bilateral $q$-series and infinite-product structure.
The remaining work is deliberately concrete: derive and verify the explicit
six-parameter Weierstrass map and translation point, implement independent
PARI/GP evaluation paths, and replace the marked placeholders with checked
numerical data.

\section*{Statement on AI assistance and computational tools}
This draft was prepared with AI-assisted synthesis, \LaTeX{} editing, and
normalization checking.  The author remains responsible for the mathematical
statements, references, computational verification, and final text.  The
planned companion implementation will use PARI/GP for exact algebra and
high-precision elliptic computations.

\begin{thebibliography}{99}
\bibitem{ScheuerleDirect2026}
T.~Scheuerle,
\newblock Direct generation of a Somos--4 sequence from an algebraic generating function,
\newblock arXiv:2609.15754 [math.CO], 2026.
\newblock \url{https://arxiv.org/abs/2609.15754}.

\bibitem{ScheuerleSix2026}
T.~Scheuerle,
\newblock A six-parameter extension of the direct Somos--4 algebraic generating function,
\newblock manuscript in preparation, 2026.

\bibitem{ScheuerleWorkflow2026}
T.~Scheuerle,
\newblock From direct algebraic generating functions to elliptic sigma and theta representations of Somos--4 sequences: A constructive workflow,
\newblock manuscript, 2026.

\bibitem{Hone2005}
A.~N.~W. Hone,
\newblock Elliptic curves and quadratic recurrence sequences,
\newblock \emph{Bull. London Math. Soc.} \textbf{37} (2005), 161--171.
\newblock \url{https://doi.org/10.1112/S0024609304004163}.

\bibitem{Hone2020}
A.~N.~W. Hone,
\newblock Continued fractions and Hankel determinants from hyperelliptic curves,
\newblock \emph{Comm. Pure Appl. Math.} \textbf{73} (2020), 2310--2347.
\newblock \url{https://doi.org/10.1002/cpa.21923}.

\bibitem{LiuWangZhang2026}
F.~Liu, Y.~Wang, and Z.~Zhang,
\newblock Hankel transform and $(\alpha,\beta)$ Somos--4 sequences,
\newblock arXiv:2608.08703 [math.CO, math.NT], 2026.
\newblock \url{https://arxiv.org/abs/2608.08703}.

\bibitem{Berndt1985}
B.~C.~Berndt,
\newblock \emph{Ramanujan's Notebooks, Part III},
\newblock Springer-Verlag, New York, 1985.

\bibitem{WhittakerWatson}
E.~T.~Whittaker and G.~N.~Watson,
\newblock \emph{A Course of Modern Analysis}, 4th ed.,
\newblock Cambridge University Press, 1927; reprinted 1996.

\bibitem{DLMF20}
NIST Digital Library of Mathematical Functions,
\newblock Chapter 20: Theta functions,
\newblock \url{https://dlmf.nist.gov/20}.

\bibitem{DLMF23}
NIST Digital Library of Mathematical Functions,
\newblock Chapter 23: Weierstrass elliptic and modular functions,
\newblock \url{https://dlmf.nist.gov/23}.

\bibitem{PARIGPManual}
The PARI Group,
\newblock \emph{User's Guide to PARI/GP},
\newblock \url{https://pari.math.u-bordeaux.fr/doc.html}.
\end{thebibliography}
\end{document}
